\section{Introduction}
Modern cloud datacenters suffer from an inherent resource mismatch: compute and memory needs rarely scale in lockstep. Many servers run out of CPU cores while gigabytes of local DRAM remain untouched, while adjacent nodes hit memory exhaustion despite largely idle processors. Paging-based memory disaggregation solves this stranded-resource problem. Compute blades keep a small working set of hot pages in local memory and access the remaining pool across dedicated memory nodes over fast interconnects like RDMA.

Operating-system-level paging makes disaggregation especially practical because unmodified binaries can execute without source code changes or specialized runtime libraries \cite{infiniswap,fastswap,hermit,dilos}. Whenever a thread accesses a missing page, the hardware traps into a page fault, and the kernel issues a 4~KiB one-sided RDMA read. On modern 100~Gb/s adapters, remote reads finish in 2--3~$\mu$s. While this is orders of magnitude faster than spinning disks or solid-state storage, it sits squarely in the ``killer microsecond'' regime described by Barroso et al.\ \cite{killer}. In such latency-critical services, tail latency (P99 and P99.9) dictates user experience far more than mean response time \cite{tail,tales}.

The core dilemma during a microsecond page fault is how to handle the CPU while waiting for the network. System builders have traditionally chosen between two extremes:
\begin{enumerate}
\item \emph{Yielding the CPU}: Stash thread registers, hand the processor to another runnable task, and wake the original thread upon packet arrival. Infiniswap pioneered this in Linux \cite{infiniswap}, but kernel context switches and interrupt processing add tens to hundreds of microseconds, completely swamping the 3~$\mu$s fetch.
\item \emph{Spinning (Busy-waiting)}: Keep the core in a tight loop polling the RDMA completion queue until the page arrives. Fastswap, Hermit, and DiLOS adopt this strategy \cite{fastswap,hermit,dilos}. While busy-waiting eliminates context-switch overhead and yields low latency when load is light, it leaves cores burning power doing zero useful computation.
\end{enumerate}

The recent EuroSys '25 paper on Adios \cite{adios} revisits this trade-off. Profiling DiLOS under realistic workloads, Yoon et al.\ exposed two fatal flaws of busy-waiting under load. First, arriving requests queue up behind spinning workers and stall for the full remote fetch time---inducing severe head-of-line (HOL) blocking that balloons tail latency. Second, because each spinning thread issues only one fetch at a time, eight workers can generate at most eight concurrent requests, leaving 100~GbE links half-utilized. Adios solves this by combining unikernel zero-privilege switching, 40-cycle user-level unithreads, and congestion-aware load balancing, achieving up to $1.64\times$ throughput and $10.89\times$ P99.9 latency cuts over DiLOS \cite{adios}.

Since building a bare-metal unikernel and RDMA testbed is outside the scope of an academic term project, we study the core question: \emph{How much of Adios's gain stems strictly from the scheduling policy once context-switch cost is fixed?} To answer this, we developed a discrete-event simulator in Python (using SimPy) that models an 8-core compute node, per-worker queue pairs, and a shared 100~GbE link. We validate our simulation against exact closed-form capacity equations across diverse hit rates, cache configurations, and context-switch latencies.

Furthermore, our profiling uncovers two key weaknesses in pure yielding: (1) under low request rates, yielding adds an unnecessary 20~ns context-switch tax while peer cores sit idle, and (2) if context switches take longer than half an RDMA fetch ($t_s > 1.5~\mu$s), pure yielding collapses below spinning throughput. To fix both problems, we implement an adaptive hybrid policy that spin-waits when cores are idle and yields immediately when queue backlog forms.

\subsection*{Summary of Contributions}
\begin{itemize}
\item \textbf{Open-source simulator}: We developed a discrete-event simulator modeling queue pairs, NIC serialization, and worker scheduling, comparing spinning and yielding under identical networking constraints.
\item \textbf{Analytical capacity bounds}: We derive closed-form formulas for saturation throughput ($\Lambda_{bw}, \Lambda_{cpu}, \Lambda_{link}$) that match our simulation measurements within 3\%.
\item \textbf{Break-even condition}: We prove and verify that yielding improves capacity if and only if two context switches cost less than one remote page fetch ($2t_s < t_f$).
\item \textbf{Comprehensive parameter sweeps}: We evaluate tail latency (P99.9), carried throughput, Zipf-skewed LRU caching, and congestion-aware dispatching against published Adios data.
\item \textbf{Adaptive hybrid policy}: We introduce a queue-aware hybrid scheduling rule that eliminates light-load yielding overhead (achieving 5.30~$\mu$s median latency) while maintaining full 2.50~MRPS saturation throughput and complete resilience to switch costs up to 1.5~$\mu$s.
\item \textbf{Methodological analysis}: We detail two subtle simulation measurement bugs (arrival-time windowing and unbounded in-flight queues) and explain how analytic bounds helped us detect and correct them.
\end{itemize}

\subsection*{Scope}
Our project is deliberately scoped to simulation. We do not run bare-metal OSv unikernels or physical Mellanox NICs. Factors like kernel memory allocators, page-table locks, and reclaimer threads are omitted to isolate the pure scheduling trade-off. Our results serve to explain the underlying queuing dynamics and identify break-even thresholds rather than benchmark real hardware.

\subsection*{Paper Outline}
The report is structured as follows. Section~\ref{sec:background} covers disaggregation fundamentals and waiting strategies. Section~\ref{sec:related} discusses related work. Section~\ref{sec:adios} reviews Adios's key mechanisms and published figures. Section~\ref{sec:model} presents our simulation model and capacity bounds. Section~\ref{sec:method} describes our experimental methodology. Section~\ref{sec:results} presents our findings, including the adaptive hybrid policy. Section~\ref{sec:analysis} uses the validated models for analytical extrapolation. Section~\ref{sec:discussion} discusses threats to validity, and Section~\ref{sec:conclusion} concludes.
